% !TeX program = xelatex
\documentclass{my_cv}
\title{Резюме}
\author{Земляная Дарья}
\date{Июнь 2024}
\usepackage[a4paper, left=2em, right=2em, top=2.5em, bottom=0em]{geometry}
\usepackage{multicol}
\setlength{\columnsep}{0cm}
\usepackage{xcolor}
\usepackage{fontspec}
\usepackage{setspace}
\usepackage[none]{hyphenat}
\usepackage{changepage}
\usepackage{hyperref}
\usepackage{graphicx}
\usepackage{tikz}
\usetikzlibrary{calc}
\usepackage{textcomp}
\usepackage{enumitem}

% Цветовая схема
\definecolor{pinkAccent}{HTML}{D63E82}
\definecolor{blueAccent}{HTML}{005F89}
\definecolor{textColor}{HTML}{333333}
\definecolor{blueAcc}{HTML}{459DBF}
\defaultfontfeatures{LetterSpace=1.5}

\hypersetup{
unicode=true,
colorlinks=true,
linkcolor=blueAccent,
urlcolor=blueAccent,
}

\setmainfont{Arial}[
BoldFont={Arial Bold},
ItalicFont={Arial Italic}
]

% Стили для элементов
\newcommand{\cvSection}[1]{
\vspace{1em}
\noindent\color{pinkAccent}\large\textbf{#1}
\vspace{0.5em}\hrule
\vspace{0.5em}
}

\newcommand{\badge}[2]{%
    \colorbox{#1}{\textbf{\ #2\ }}%
}

\newcommand{\positionAndCompanyLine}[4]{
\vspace{-0.3em}
\noindent\textbf{#1} \
\noindent #2 \hfill #3 -- #4 \
\vspace{0.2em}
}

% Добавляем определение techStackList
\newcommand{\techStackList}[1]{
\vspace{0.1em}
\noindent\textit{Tech Stack: #1}
\vspace{0.3em}
}

\setlist{nosep, leftmargin=1.2em}

\begin{document}

\begin{multicols}{3}
    \headerSection{\badge{pinkAccent}{Системный программист}}
    \columnbreak
    \headerSection{\xhrulefill{pinkAccent}{0.6em}}
    \columnbreak
    \headerSection{\xhrulefill{blueAccent}{0.6em}}
\end{multicols}

\begin{multicols}{2}
\begin{minipage}[t]{1\linewidth}
\vspace{-2em}
\headerSection{\Huge \textbf{\textcolor{textColor}{Земляная Дарья}}}
\cvParagraph{%
\href{tel:+79091559499}{+7(909) 155 94 99}\

\href{mailto:darijazeml@inbox.ru}{darijazeml@inbox.ru}\

\href{https://github.com/Dariazeml1007}{github.com/Dariazeml1007}
}

   \vspace{0.5em}
\noindent
\begin{minipage}[t]{0.65\linewidth}
    \textbf{{Основные навыки:}}\\[\smallskipamount]
    \textcolor{textColor}{C \textbullet\ C++ \textbullet\ x86/x64 ASM \textbullet\ Python \textbullet\ Ruby}\\[\medskipamount]
    \textbf{{Инструменты:}}\\[\smallskipamount]
    \textcolor{textColor}{Git \textbullet\ CMake \textbullet\ GDB \textbullet\ Valgrind \textbullet\ Linux}\\[\medskipamount]
    \textbf{{Языки:}}\\[\smallskipamount]
    \textcolor{textColor}{Русский (родной) \textbullet\ Английский (B1)}
\end{minipage}
\end{minipage}
\hfill
\begin{minipage}[t]{0.8\linewidth}
    \vspace{-1.5em}
    \begin{flushright}
        \begin{tikzpicture}
            \node[draw=pinkAccent, line width=1.8pt,
                  double=blueAccent, double distance=1.2pt,
                  rounded corners=4pt,
                  inner sep=0pt] (frame) {
                \includegraphics[width=4.8cm,height=6.5cm,keepaspectratio]{photo.jpg}
            };
        \end{tikzpicture}
    \end{flushright}
\end{minipage}
\end{multicols}

\setlength{\columnseprule}{0pt}
\setlength{\columnsep}{2em}

% === ОБРАЗОВАНИЕ ===
\cvSection{\badge{blueAcc}{Образование}}
\begin{adjustwidth}{2em}{0em}
\positionAndCompanyLine{Московский Физико-Технический Институт}{}{2024 -- н.в.}
\positionAndCompanyLine{Лицей "Вторая школа"}{}{2021 -- 2024}
\end{adjustwidth}

% === СПЕЦКУРСЫ ===
\cvSection{\badge{pinkAccent}{Спецкурсы}}
\begin{adjustwidth}{2em}{0em}
\begin{multicols}{2}
  \begin{itemize}
    \item \textbf{Компиляторные технологии на C} (Дединский И.Р)
    \item \textbf{Промышленное программирование на C++} (Владимиров К.И)
    \item \textbf{Архитектура вычислительных систем} (Sber)
    \item \textbf{Тензорные компиляторы} (Sber)
    \item \textbf{Тестовые генераторы}
  \end{itemize}
\end{multicols}
\end{adjustwidth}

% === ОПЫТ РАБОТЫ ===
\cvSection{\badge{blueAcc}{Опыт работы}}
\begin{adjustwidth}{2em}{0em}
\positionAndCompanyLine{Стажер по анализу данных}{МТС Digital}{06/2024 -- 08/2024}
\begin{flushleft}\textbf{\textit{Big Data стажировка}} \end{flushleft}
\begin{itemize}
\item Работа с PySpark для обработки больших данных
\item Использование Git/GitLab в командной разработке
\end{itemize}
\techStackList{PySpark, Git, Python}
\end{adjustwidth}

% === ПРОЕКТЫ ДЕДИНСКИЙ ===
\cvSection{\badge{pinkAccent}{Проекты @ Huawei/МФТИ (Дединский И.Р.)}}
\begin{adjustwidth}{2em}{0em}

\positionAndCompanyLine{Эмулятор процессора}{C}{2024}
\begin{itemize}
\item Ассемблерный компилятор + стековая ВМ на чистом C
\item Поддержка меток, примеры рекурсии (факториал)
\item \href{https://github.com/Dariazeml1007/Processor}{[Ссылка на github]}
\end{itemize}
\techStackList{C}

\positionAndCompanyLine{Низкоуровневая реализация printf}{C, Nasm}{2025}
\begin{itemize}
\item Интерфейс C-ASM через extern (System V ABI)
\item Поддержка \%d, \%s, \%c, \%\%, \%b/\%o/\%x
\item \href{https://github.com/Dariazeml1007/MyAsmPrintf}{[Ссылка на github]}
\end{itemize}
\techStackList{Nasm, System V ABI}

\positionAndCompanyLine{Оптимизация множества Мандельброта}{C, AVX2}{2025}
\begin{itemize}
\item Сравнение скалярной, -O3 и ручной AVX2 реализаций
\item Интерактивная визуализация с TxLib
\item \href{https://github.com/Dariazeml1007/Mandelbrot}{[Ссылка на github]}
\end{itemize}
\techStackList{C, AVX2 Intrinsics}

\positionAndCompanyLine{Оптимизация хэш-таблицы}{C, AVX2, ASM}{2025}
\begin{itemize}
\item Оптимизация через CRC32/AVX2, профилирование Valgrind
\item Ускорение поиска в 1.76× (коэффициент 220)
\item \href{https://github.com/Dariazeml1007/Hash_table}{[Ссылка на github]}
\end{itemize}
\techStackList{C, x86-64 ASM, AVX2, Valgrind}

\end{adjustwidth}

% === ПРОЕКТЫ ВЛАДИМИРОВ ===
\cvSection{\badge{blueAcc}{Проекты по C++ (Владимиров К.И.)}}
\begin{adjustwidth}{2em}{0em}

\positionAndCompanyLine{2Q Cache Simulator}{C++17}{2025}
\begin{itemize}
\item Реализация Two-Queue алгоритма, точность ~90\%
\item \href{https://github.com/Dariazeml1007/Cache_2Q_simulation}{[Ссылка на github]}
\end{itemize}
\techStackList{C++17, STL}

\positionAndCompanyLine{BST с range queries}{C++17}{2025}
\begin{itemize}
\item Подсчёт элементов в диапазоне, балансировка, dump в Graphviz
\item \href{https://github.com/Dariazeml1007/BST-Range-Query}{[Ссылка на github]}
\end{itemize}
\techStackList{C++17, AVL Trees, GoogleTest}

\positionAndCompanyLine{Интерпретатор ParaCL}{C++17, Flex/Bison}{2026}
\begin{itemize}
\item Лексер (Flex) → парсер (Bison) → AST → интерпретация
\item Модульные тесты GoogleTest
\item \href{https://github.com/Dariazeml1007/ParaCL}{[Ссылка на github]}
\end{itemize}
\techStackList{C++17, Flex, Bison, GoogleTest}

\end{adjustwidth}

% === ПРОЕКТЫ СИМУЛЯЦИЯ ===
\cvSection{\badge{pinkAccent}{Задача по курсу "Средства симуляции"}}
\begin{adjustwidth}{2em}{0em}

\positionAndCompanyLine{Эмулятор CPU c TOY ISA}{C++, Ruby}{2025}
\begin{itemize}
\item 14 инструкций, mixed encoding, поддержка прерываний
\item Микро-ассемблер на Ruby → эмуляция на C++
\item \href{https://github.com/Dariazeml1007/CPU-interpreter}{[Ссылка на github]}
\end{itemize}
\techStackList{C++17, Ruby, DSL}

\end{adjustwidth}

% Небольшой отступ для красивого завершения
\vspace{0.3cm}

\end{document}
